\section{Essence of Archaism}

In \textit{The Essence of Archaism}\footnote{\url{http://guillaumefayearchive.wordpress.com/2007/07/11/the-essence-of-archaism}}, Guillaume Faye describes the return of what he calls ``archaic values''.

First of all, the Greek word \textit{arche}, like the Latin \textit{principium}, means what comes first, either temporarily (``beginning'') or ontologically (``principle''). Thus the Prologue to John's Gospel in Greek is, ``En arche en o logos'', and in Latin, ``in principio erat verbum'' (``In the Beginning was the Logos''). This we can understand to mean, ``The foundation of order is the first principle.'' This is the foundation of Christendom, the ancient pagan civilizations, the Bhagavad Gita according to B G Tilak\footnote{\url{https://www.gornahoor.net/?p=1619}}, and the Tao Te Ching as Matgioi informs us.

Mr. Faye makes a list of archaic values, which are listed here:
\begin{itemize}
\item distinctive sexual roles
\item    the transmission of ethnic and popular traditions
\item    spirituality and sacerdotal organization
\item    visible and supervisory social hierarchies
\item    the worship of ancestors
\item    initiatory rites and tests
\item    the reconstruction of organic communities that extend from the individual family unit to the overarching national community of the people
\item    the deindividualization of marriage to involve the community as much as the couple
\item    the end of the confusion of eroticism and conjugality
\item    the prestige of the warrior caste
\item    social inequality, not implicit, which is unjust and frustrating, as in today’s egalitarian utopias, but explicit and ideologically justifiable
\item    a proportioned balance of duties and rights
\item    a rigorous justice whose dictates are applied strictly to acts and not to individual men, which will encourage a sense of responsibility in the latter
\item    a definition of the people and of any constituted social body as a diachronic community of shared destiny, not as a synchronic mass of individual atoms
\end{itemize}

He attributes this list to \textbf{Raymond Ruyer}, whose books are all apparently out of print and difficult to find. It seems to be a good list, most of which can also be found scattered around Gornahoor. The difficulty is determining the source of those values.

Mr. Faye suggests they are biological values, which seem untenable (if the phrase is even coherent), since biological laws (not values) are immutable. It would be surprising to find beavers building dams out of poured concrete, or honey bees to suddenly build their societies around a King bee. But human societies that abrogate all those those archaic values do exist, so a biological explanation seems untenable.

So we are left with human values. Obviously, we need to know what a human being is, before we can suppose those archaic values to specifically human. Then the question arises about why humans would choose disvalue instead of value. That Gornahoor has suggested is due to the cosmic cycle, the degeneration of castes, and the incarnation of inferior beings.

So we are left with the true source: the values are transcendental, as the very term ``archaic'' implies. Those archaic values follow the principle of order; to choose otherwise is to be in rebellion against the Logos. That is the essence of revolution.

As for the return to such values, either deliberately or from their imposition by Islam, Rene Guenon\footnote{\url{https://www.gornahoor.net/index.php?s=prophecies+rene+guenon\&submit=Search}} has already shown us the options:

\begin{itemize}
\item Degeneration
\item Assimilation
\item Transformation
\end{itemize}


As for the press baron complaining about the problems of capitalism, which he expects Islam to eliminate: Mr. Faye should have looked to the recent past of his country, particularly the works of \textbf{Charles Maurras}, like him, a pagan with an appreciation for the civilization of Christendom.

One of the key points of Action Française was Catholic social teaching. Since it has its foundations in natural law and Roman jurisprudence, it is universal, and not solely Catholic. It addresses issues such as just prices and wages, competition, the problem of usury. Neither socialist nor capitalist, it proposes a widespread private ownership of the means of production, with competition based on quality and reliability. Through the prohibition of usury, those with access to capital would not be able to simply buy up the market, and thereby, effectively eliminating competition.

That social teaching was not ``imposed'' on Europe, it was accepted as archaic value. The real question that Mr. Faye does not answer, since he never even asks it, is ``who imposed the alternative?''


